\section{政策含义}
\label{sec:policy}

前述结果对算法审计、资产管理、风险管理与投资者保护各有一层含义。
以下五项均以本文估计的具体量为基础，指向可操作的判据或参数。

算法审计的三个诊断量。定理 \ref{thm:threshold} 把门控占优的条件化为三个量：互补性统计量 $\hat\Gamma$、
噪声底线 $\Theta$ 与成本项 $2\lambda\kappa\mathcal{C}(\hat g)$。
三者都可以由第三方在同一数据集上重算：$\hat\Gamma$ 由制度内的条件信息系数
或专家组合收益估计；$\Theta$ 由区块自助法在无真实互补性的原假设下自校准，
不依赖对外生参数的假设；成本项由换手率与费率参数直接计算。
据此，智能投顾与量化资管产品的算法备案与信息披露可以引入三个诊断量。

主动暴露为组合权重相对基准的半 $L_1$ 距离均值，作用是识别机制退化。
在绝对收益目标下主动暴露降至 $0.0003$，此时任何自适应宣称都不成立；
该指标只依赖持仓序列，计算成本极低。
互补性统计量与噪声底线的对比用于区分机制存在与机制可识别。
本文中该对比的 $p$ 值为 0.060 与 0.065，机制存在但不可稳健识别，
这一区分直接决定投资者对超额收益宣称的合理预期。
一阶与二阶渠道的分别检验用于识别收益来源。本文证据显示门控收益
来自二阶风险配置渠道，监管评估的重点相应落在风险暴露的真实性与稳定性上，
而非预测能力的宣称。三个诊断量不涉及重新训练模型，
足以剔除退化均衡、抽样噪声与来源误判三类问题，
与监管科技强调的可验证、可复现原则一致\cite{chakraborty2017,bis2024aer}。

资产管理中的容量与成本匹配。第 \ref{sec:marginal} 节的估计显示，提高模型复杂度的边际回报很小甚至为负。
单边交易成本每上升 10 个基点，门控的相对优势下降约 0.015，盈亏平衡成本约 92 个基点；
制度状态数由 3 增至 4 使相对优势下降 0.014；训练窗长度没有单调效应
（斜率 $-0.006\,3$/年，$p=0.617$）。三项估计指向三条管理规则。

换手控制优先于模型精细化。门控组合的换手率为 0.372，约为等权基准的 2.8 倍；
即使在零成本假设下，其相对优势也仅 $+0.135$。
成本不是限制机制的根本原因，却是净价值最敏感的维度。
制度数与样本量必须匹配。$T=247$ 的样本量下，200 个自由参数的门控处于容量约束的
紧边界上；本文实证最优的制度数为 3，继续细分损害净价值。
训练窗并非越长越好。3 年训练窗的相对优势（$+0.213$）高于 8 年（$+0.163$）
与 5 年（$+0.120$）。在市场结构快速变化的环境下，扩大训练样本会引入陈旧信息，
模型治理应把数据新鲜度与容量约束一并评估。

风险管理中的基线选择。渠道分解表明门控收益主要来自对不同制度下风险结构的适应，而非对收益水平的预测。
这一来源有两个含义。收益的稳定性较高：风险结构（波动率与相关性）的制度差异
比收益水平的制度差异更稳定，二阶渠道的价值在时间上更可持续。
复杂架构的必要性下降：低波动、风险平价这类纯化的风险结构策略是更高效的实现方式。
表 \ref{tab:main} 与此一致，低波动专家（夏普 0.804）与风险平价（0.769）
表现优于包含完整门控机制的组合（0.673），并在 5\% 水平上稳健显著
（Holm $p=0.041$ 与 $0.017$）。

风险管理部门可以据此设定配置流程：先用简单风险结构策略建立基线，
只有当自适应机制在扣除成本后持续超越该基线，才承担其模型风险与运维复杂度。
尾部风险维度提供了另一个取舍参照。门控相对等权的日频 5\% CVaR 改善为 $+0.0047$
（区间含零），换手率提高 $+0.241$（Holm $p=0.005$）。
以控制 CVaR 为目标的机构，直接施加尾部约束或采用风险平价，
比依靠门控换取尾部改善更为有效。

投资者保护中的披露边界。三类稳健性证据界定了业绩宣称的合理范围。时间上，门控相对优势由 2014--2019 年的
$+0.180$ 降至 2020--2025 年的 $+0.049$，全样本平均掩盖了这一衰减，
披露应包含分时段业绩。设定上，四种制度定义下的夏普跨度约 0.16，
与门控的整体优势（$+0.12$）同量级，制度变量选择这一研究者自由度
本身就贡献了与结论同量级的不确定性，披露应包含关键设定的敏感性区间。
推断上，阈值检验在区块长度 3--6 个月时显著、12--24 个月时不显著，
统计显著性处于视推断方法而定的边缘地带。

智能投顾产品的业绩说明可以据此要求三项内容：分时段业绩与子样本稳定性、
关键设定（制度定义、制度数、训练窗）的敏感性区间、
以及预先注册的推断方法（配对单位、自助区块长度、多重比较校正方式）。
这些要求对依赖短样本、单一随机种子或模拟市场的评价尤为必要：
本文的证据显示，同一策略在同一市场上，仅改变推断方法就足以在显著与不显著之间移动结论。

市场效率与系统性风险。本文的结果对市场效率与金融稳定各有一层含义。

对市场效率的含义。门控价值的存在说明状态信息尚未被完全定价：
若市场在半强式意义上有效，按状态调整信号权重不应带来稳定收益\cite{grossman1980}。
本文估计的门控优势（年化主动收益 3.77\%、夏普差值 $+0.115$）在统计上与零不可区分，
且盈亏平衡成本约为 92 个基点，与 A 股主动权益投资的显性费率量级接近。
两项事实合起来的读法是：状态信息的价格已接近其获取成本，
市场在该维度上接近有效边界，但尚未越过。这一判断可以随算法能力提升而移动，
其移动速度取决于估计误差成本 $\Theta$ 能否随数据频率与横截面宽度下降（定理 \ref{thm:channels}）。

对系统性风险的含义。制度状态的高度持续性（半衰期 3.25 个月）
与约 19.5 倍的持续性乘数，意味着状态识别在月度频率上可行，
风险管理部门有条件在制度切换早期调整暴露。需要关注的是同质化：
当大量机构采用结构相似的多智能体框架、依据同一类公开状态变量调整权重时，
制度切换可能成为一致行动的触发点。市场流动性与资金流动性之间存在相互强化的机制，
交易者的融资能力取决于其持仓的市场流动性，市场流动性又取决于交易者可用的资金，
两者在压力情景下形成流动性螺旋\cite{brunnermeier2009}。
状态自适应策略在切换点上集中减仓，正是这一机制的现实触发路径。
本文的年度环境证据与此一致：门控价值集中于高波动、高制度切换的年份
（2020 年制度切换 7 次，差值为 $+0.305$；2021 年 6 次，差值为 $+0.355$），
这类年份同时是拥挤交易与流动性风险容易累积的时段。
冲击年证据补充了一个细节：门控价值与状态分化的幅度同向，
与切换次数关系不大（表 \ref{tab:shock}）。策略拥挤的风险因此主要不在
「切换有多频繁」，而在「多少资金在相似的幅度阈值上行动」。

前瞻性的制度设计。基于上述判断，本文提出三项可以落地为规则的设计。
其一，状态变量与切换阈值的披露：算法备案除要求披露模型类型外，
可要求披露所使用的状态变量清单与调仓触发条件，
使监管者能够识别同类策略的重叠程度。
其二，压力测试情景纳入状态切换：现有压力测试多以价格冲击为情景，
可增加「状态变量集中触发」情景，即假定多数同类策略在同一阈值附近调仓，
评估由此产生的流动性与波动率联动。
其三，拥挤度监测指标：本文的互补性统计量 $\hat\Gamma$、主动暴露与换手率
可以直接作为监测指标——$\hat\Gamma$ 反映状态信息被利用的程度，
主动暴露反映策略偏离基准的幅度，换手率反映调仓强度。
三者均为持仓层面可计算的量，不依赖策略内部模型，
可用于估算采用状态自适应框架的资金规模及其集中度。
这一设计的意义在于把监管从「事后评估业绩」推进到「事前识别共振条件」。

上述五方面含义可归纳为三点。可复核性优先：主动暴露、互补性与噪声底线、
渠道分别检验三个低成本诊断量，把智能投顾的业绩宣称转化为可审计的对象。
容量与成本的匹配优先于模型复杂度：制度数与训练窗的边际回报有限甚至为负，
成本是净价值最敏感的维度。简单策略作为默认基线：门控收益主要来自风险配置，
低波动与风险平价在本文样本上更为有效，模型风险与运维成本也低于多智能体系统。
